\Article{32}
Wie groß ist aber der Winkel $A$, wenn die Anziehung des Gefäßes größer ist als die gegenseitige Anziehung der Flüssigkeitsteile? Der imaginäre Wert, den für $\beta\beta>\alpha\alpha$ die Formel $\sin\frac12 A=\frac\beta\alpha$ dem Winkel $A$ zuweist, zeigt bereits, dass eine der zugrunde liegenden Annahmen in diesem Fall unzulässig ist. Tatsächlich lässt sich, wenn $\beta\beta>\alpha\alpha$, die Annahme einer \textit{Begrenzung} der Oberfläche $T$ mit der Bedingung eines Minimums der Funktion $W$ nicht vereinbaren. Wo man die Grenze auch ansetzt, denke man sich jenseits davon einen sehr dünnen Flüssigkeitsfilm ausgebreitet, sodass $T$ um $T'$ und demnach $U$ um annähernd ebenso viel zunimmt. Der Wert der Funktion $W$ ändert sich dann offenbar um einen Betrag, der der negativen Größe $-(2\beta\beta-2\alpha\alpha)T'$ bis auf einen unmerklichen Unterschied gleich ist. Der Wert von $W$ kann sogar weiter abnehmen, bis $T'$ die ganze übrige Gefäßoberfläche bedeckt. Der Wert $-(2\beta\beta-2\alpha\alpha)T'$ für die Änderung ist umso genauer, je kleiner man die Dicke annimmt. Solange es nur um den Wert des Ausdrucks $W$ geht, spricht nichts dagegen, sich die Dicke bis zum Verschwinden vermindert zu denken. Ein Film von verschwindender Dicke, die man von einer unmerklichen Dicke sorgfältig unterscheiden muss, wäre jedoch eine bloße mathematische Fiktion. Die dieser Fiktion entsprechende Gestalt des Raumbereichs $s$ wäre tatsächlich dieselbe wie diejenige, für die $W$ im Fall $\beta\beta=\alpha\alpha$ einen kleinsten Wert annimmt.

In unserer physikalischen Aufgabe verhält es sich etwas anders. Damit Gleichgewicht möglich ist, muss ein solcher zusätzlicher Film eine bestimmte, wenn auch unmerkliche Dicke haben. Ist ein solcher Teil vorhanden, so ist der Ausdruck $W$ unvollständig, wie wir in Artikel 18 gezeigt haben. Bezeichnen wir den vom Film bedeckten Teil des Gefäßes mit $T'$
\SourceLeaf{76}und seine Dicke in einem beliebigen Punkt mit $\rho$, so sind zum Ausdruck $\Omega$ noch die folgenden Terme hinzuzufügen
\[
\pi cc\int\theta'\rho.\mathrm{d}T'-\pi cC\int\Theta'\rho.\mathrm{d}T'
\]
und damit zum Wert von $W$ diese
\[
\begin{gathered}
\frac{\pi C}g\int\Theta'\rho.\mathrm{d}T'-\frac{\pi c}g\int\theta'\rho.\mathrm{d}T'\\
=\int\mathrm{d}T'.\left(\frac{2\beta\beta}{\Theta0}.\Theta'\rho-\frac{2\alpha\alpha}{\theta0}.\theta'\rho\right)
\end{gathered}
\]
Da also der Wert von $W$ durch das Hinzukommen dieses Films bereits die Änderung $(2\beta\beta-2\alpha\alpha)T'$ erfahren hat, beträgt die gesamte Änderung, die zu dem ohne den Film geltenden Wert von $W$ hinzuzufügen ist,
\[
-2\int\mathrm{d}T'.\left[\beta\beta\left(1-\frac{\Theta'\rho}{\Theta0}\right)-\alpha\alpha\left(1-\frac{\theta'\rho}{\theta0}\right)\right]
\]
Wegen $\theta'0=\theta0$ und $\Theta'0=\Theta0$ wäre diese Änderung bei verschwindender Dicke null. Da aber $\theta'\rho$ und $\Theta'\rho$ bei wachsender Dicke $\rho$ sehr rasch abnehmen und bereits für einen unmerklichen Wert von $\rho$ unmerklich werden, nähert sich diese Änderung sehr rasch dem Wert $-(2\beta\beta-2\alpha\alpha)T'$. Im Gleichgewichtszustand der Flüssigkeit muss sie diesem Wert bis auf einen unmerklichen Unterschied gleich sein, damit der Wert des korrigierten Ausdrucks $W$ keine weitere merkliche Abnahme erfahren kann. Eine vollständige Untersuchung des Gesetzes, dem die Dicke $\rho$ folgen muss, würde eine eingehendere Rechnung erfordern. Dabei verweilen wir hier jedoch nicht, weil sie ohne Kenntnis der Funktionen $f,F$, von denen die Funktionen $\theta',\Theta'$ abhängen, und aus den in Artikel 34 anzugebenden Gründen wenig ergiebig erscheinen könnte. Zur Untersuchung des Hauptteils der Flüssigkeit, also des Teils, dessen Abmessungen sämtlich merklich sind, genügt es in unserem Fall mit $\beta\beta>\alpha\alpha$, sich das Gefäß in der Nähe der Grenze des Hauptteils \textit{benetzt} zu denken. Es ist also von einem Flüssigkeitsfilm bedeckt, dessen Dicke zwar unmerklich ist, aber groß genug, dass $\theta'\rho$ und $\Theta'\rho$ vernachlässigt werden können. Die Funktion, die im Gleichgewichtszustand ein Minimum sein muss, lautet dann
\[
\int z\mathrm{d}s-2(\beta\beta-\alpha\alpha)(T+T')-\alpha\alpha T+\alpha\alpha U
\]
wobei $T,U$ nur auf den Hauptteil der Flüssigkeit bezogen werden. Offenbar stimmt also die Variation dieser Funktion infolge einer virtuellen Änderung der Gestalt des Hauptteils der Flüssigkeit, bei der die Summe $T+T'$ unverändert bleibt, mit der Variation des Ausdrucks
\SourceLeaf{77}%
\[
\int z\mathrm{d}s-\alpha\alpha T+\alpha\alpha U
\]
überein, also genau des Ausdrucks, der im Fall $\beta\beta=\alpha\alpha$ ein Minimum sein muss. Daraus schließen wir, dass die Gleichgewichtsform der Flüssigkeit in einem Gefäß, für das $\beta\beta>\alpha\alpha$ gilt, mit der Gleichgewichtsform derselben Flüssigkeit in einem Gefäß übereinstimmt, für das $\beta\beta=\alpha\alpha$ gilt. Der Unterschied besteht darin, dass sie im ersten Fall im strengen Gleichgewicht in einen Film von unmerklicher Dicke auslaufen muss. Der berühmte \textsc{Laplace} hat übrigens bereits darauf hingewiesen, dass im ersten Fall ein mit einem Flüssigkeitsfilm von unmerklicher Dicke bedecktes Gefäß einem Gefäß gleichwertig ist, dessen Teilchen auf die Flüssigkeit dieselbe Anziehungskraft ausüben wie die Flüssigkeitsteile aufeinander.

Daraus ergibt sich unmittelbar die Änderung, die an den Aussagen in Artikel 18 über den Anstieg von Flüssigkeiten in senkrechten Kapillarröhren vorzunehmen ist. Gilt $\beta\beta>\alpha\alpha$, so ist in den dort angegebenen Formeln $\alpha\alpha$ anstelle von $\beta\beta$ einzusetzen.

\Article{33}
Im Fall $\beta\beta<\alpha\alpha$ kann das Gefäß nicht von einem Flüssigkeitsfilm unmerklicher Dicke benetzt werden, wenn das Gesetz der Funktionen $\theta',\Theta'$ so beschaffen ist, dass der Wert der Funktion
\[
\alpha\alpha\left(1-\frac{\theta'\rho}{\theta'0}\right)-\beta\beta\left(1-\frac{\Theta'\rho}{\Theta'0}\right)
\]
für die wir der Kürze halber $Q\rho$ schreiben, stetig zunimmt, wenn $\rho$ vom Wert $0$ zu einem merklichen Wert fortschreitet. Bei diesem Verhalten der Funktion $Q\rho$ würde das Vorhandensein eines solchen Films offenbar der Bedingung des Minimums widersprechen. Dieses Verhalten ergibt sich unmittelbar aus der in Artikel 31 besprochenen Annahme, nach der $fx,Fx$ in einem bestimmten, von $x$ unabhängigen Verhältnis stehen. Daraus folgt nämlich auch $\frac{\theta'\rho}{\theta'0}=\frac{\Theta'\rho}{\Theta'0}$ und somit $Q\rho=(\alpha\alpha-\beta\beta)(1-\frac{\theta'\rho}{\theta'0})$. Folgten aber die Funktionen $f,F$ verschiedenen Gesetzen, so wäre es nicht unmöglich, dass der Wert von $\frac{\theta'\rho}{\theta'0}$ rascher abnähme als der von $\frac{\Theta'\rho}{\Theta'0}$ und die Funktion $Q\rho$ im Bereich der unmerklichen Werte von $\rho$ zunächst negativ würde und dann, nachdem sie ihren kleinsten Wert, also ihr negatives Extrem, erreicht hätte, durch den Wert $0$ wieder bis zu ihrer positiven Grenze $\alpha\alpha-\beta\beta$ anstiege. In einem solchen Fall würde das Gleichgewicht tatsächlich einen unmerklichen Film erfordern, dessen Dicke im Allgemeinen so groß sein müsste, dass $Q\rho$ sich von seinem kleinsten Wert nicht merklich unterscheidet. Bezeichnen wir diesen Wert mit $-\beta'\beta'$, so gilt $\beta'\beta'<\beta\beta$. Die Gestalt des Hauptteils der Flüssigkeit bestimmt sich dann genauso wie in einem Gefäß, für das anstelle der Größe $\beta\beta$ die Größe $\beta'\beta'$ einzusetzen wäre. Das heißt, der Winkel zwischen der Ebene,
\SourceLeaf{78}die die freie Flüssigkeitsoberfläche an der Grenze des Hauptteils berührt, und der Gefäßwand ist $=2\operatorname{arc.}\sin\frac{\beta'}{\alpha'}$. Da es aber sehr zweifelhaft ist, ob ein solcher Fall in der Natur vorkommt, scheint es überflüssig, länger dabei zu verweilen.

